\chapter[SYSTEM IMPLEMENTATION AND METHODOLOGY]{System Implementation and Methodology}

This chapter traces the process of implementing from generated loan records to the assessment returned by the API. It provides an explanation of pre-processing, financial features, training of XGBoost, SHAP and LIME, and the actual validation formulas used by the program. The last sections are dedicated to the interface, persistence and PDF workflow.

\section{Development Setup}
\subsection{Technology Stack}
The backend is written in Python with FastAPI, Pydantic, SQLAlchemy, pandas, NumPy, scikit-learn, XGBoost, SHAP, LIME and ReportLab. Frontend is a React app with Vite and Tailwind CSS. Axios is used to make API calls. There is one loan-prediction router, and a registry that loads the home-loan, personal-loan, and legacy model pipelines. The main product paths are the home-loan and personal-loan pipelines, the legacy pipeline is for backwards compatibility.

\subsection{Development Environment}
The repository README provides information on Python 3.10+ and Node.js 18+ and separate backend and frontend development servers. It is usually served by Uvicorn on port 8000 and Vite serves the client during development. Dependencies are listed in \texttt{requirements.txt} and \texttt{frontend/package.json}. The Python requirements do not specify a fully frozen environment, and exact reproducibility also requires package versions and platform.

\section{Data Preparation and Model Development}
\subsection{Dataset Description and Inspection}
There are two datasets of products: \path{datasets/homeloan_train.csv} and \path{datasets/personalloan_train.csv}. Each of the repository scripts creates 3,000 records with NumPy random seed 42. There are 18 columns in the home-loan table and 17 columns in the personal-loan table. They both have a loan ID, attributes of the applicant, financial inputs, credit information, and a label called \texttt{Loan\_Status}. There are 968 approved and 2,032 rejected home-loan records in the data set. The personal-loan data set contains 1,368 approved loans and 1,632 rejected loans. The checked CSVs do not contain any empty cells or duplicate full rows.

These are data created by the repository that are not real. The label-generation scripts use hand-specified financial and credit rules without a gender-specific adjustment. Gender remains a model input and audit field. The high test scores reported later are therefore scores on a designed simulation, rather than scores on the applicant population of a bank's past or future. There is no externally sourced lending dataset or data dictionary in the repository.

Graph~\ref{graph:dataset-classes} shows the observed approval and rejection counts for the two products.

The class proportions differ between products, with more rejections in both datasets. They do not set a realistic approval rate for a lending institution. The proportions are determined by the distributions and decision rules selected. Those correlations introduced by the generators (e.g., education and income) are part of that authored setting.

\reportgraph{figures/graphs/dataset_class_distribution.pdf}{0.90}{6.1}{Generated dataset class counts}{Class counts read from the current home-loan and personal-loan CSVs.}{graph:dataset-classes}

\subsection{Data Preprocessing}
The training scripts read in the product CSV, calculate derived features, extract \texttt{Loan\_Status} from the predictors, and remove \texttt{Loan\_ID} from the predictors. The target is assigned the value of 1 if it is Approved and 0 if it is Rejected. Random state 42 is used to split the train into an 80/20 stratified split. The preprocessor is a column transformer, which is a pipeline of transformations: categorical fields are transformed by one-hot encoding, where the first category is dropped and unknown values are ignored; numerical fields are transformed by standard scaling. The transformer is fit on the training set and then applied to the test set.

If you make a request online, Pydantic schemas will normalize optional numeric values, while the model registry will provide defaults for the product. This normalization is convenient for operations, but it may be possible to fill in default values for values that are not given; the user should enter valid application data and check the form of data that is displayed as submitted. Batch processing normalizes rows, and records rows that are not valid as flawed records that include an explanation of the validation problem.

\subsection{Financial Feature Engineering}
The two products use distinct financial calculations. For a loan principal $P$, annual interest rate $i$, and term $n$ months, let $r=i/12$. The amortized monthly instalment used by the product pipelines is
\begin{equation}
  \mathrm{EMI}=P\frac{r(1+r)^n}{(1+r)^n-1}.
  \label{eq:emi}
\end{equation}
The home-loan generator and preprocessing function have an illustrative annual rate of 8.5\%, the personal-loan path has a rate of 13.5\%. These are implementation assumptions and not real lender quotes.

Monthly household income is the income of the applicant plus co-applicant for a home loan. The FOIR denominator for a personal loan is the income of the applicant. In both product-specific pipelines the ratio is
\begin{equation}
  \mathrm{FOIR}=\frac{\mathrm{Existing\ EMI}+\mathrm{Proposed\ EMI}}{\mathrm{Income\ used\ for\ product}}.
  \label{eq:foir}
\end{equation}
The home-loan pipeline also calculates
\begin{equation}
  \mathrm{LTV\ Ratio}=100\times\frac{\mathrm{Loan\ Amount}}{\mathrm{Property\ Value}}.
  \label{eq:ltv}
\end{equation}
Personal-loan features include loan purpose and credit-card utilization instead of property collateral features. The raw fields for home-loans are property value and property area; for personal-loans, the raw fields are purpose and utilization. Each model has a numerical feature set that includes engineered EMI/FOIR and product-specific columns.

\paragraph*{Worked Financial Calculations}
Consider a home-loan request for Rs.~3,000,000 over 240 months. At the feature function's annual rate of 8.5\%, the proposed monthly instalment is approximately Rs.~26,034.70. With household income of Rs.~80,000 and existing monthly repayments of Rs.~10,000, the FOIR is
\[
\frac{26,034.70+10,000}{80,000}\approx 0.4504.
\]
If the property value is Rs.~4,000,000, the LTV ratio is 75\%. These are feature calculations, not an example of a measured approval prediction. The full profile, including credit and categorical fields, is still sent to the classifier.

If you take a personal loan for Rs.~300,000 for 36 months, the feature rate is 13.5\%, which means that your monthly instalment will be approximately Rs.~10,180.59. The FOIR is approximately 30.36\%, with the applicant income being Rs.~50,000 and the existing repayments being Rs.~5,000. This product's model FOIR does not include co-applicant income. This difference is one of the reasons why the feature functions are product-specific.

The API also creates a solvency summary. If provided, its general income display will show the co-applicant income, and the personal-loan feature function will use applicant income. Users are therefore advised to read the named product feature and displayed summary in their own formula and not to take it for granted that all income fields are performing the same function.

\subsection{XGBoost Training and Evaluation}
The XGBoost classifiers are trained separately in \path{scripts/train_homeloan.py} and \path{scripts/train_personalloan.py}. The two scripts share the same number of estimators (120), maximum tree depth (4), learning rate (0.08), row subsampling (0.85), column subsampling (0.85), random state (42), and evaluation metric (log-loss). Each script saves the classifier, fitted preprocessor, transformed feature names, a 200-row transformed background sample, and metrics under the product directory. Those artifacts are loaded into a \texttt{ProductPipeline} during runtime with the product feature function and explainers.

\noindent\begin{minipage}{\textwidth}
\textbf{Algorithm 1: Product-specific model training}\par\smallskip
\begin{lstlisting}
Input: Product dataset D, product feature function g, seed 42
Output: Fitted preprocessor T, XGBoost classifier M, test metrics
1. D' <- g(D); remove Loan_ID from predictors
2. Encode Loan_Status: Approved=1, Rejected=0
3. Split predictors and labels with stratification (80% / 20%)
4. Fit one-hot encoding and standard scaling on training predictors
5. Transform training and test predictors with the fitted transformer
6. Fit XGBoost(n_estimators=120, max_depth=4, learning_rate=0.08,
   subsample=0.85, colsample_bytree=0.85, random_state=42)
7. Predict labels and approval probabilities on the held-out test split
8. Save model, transformer, feature names, background sample, and metrics
\end{lstlisting}
\end{minipage}\par\medskip

The training scripts compute accuracy, precision, recall, F1-score and ROC-AUC on the held-out split. Product-specific metrics are saved in \texttt{models/homeloan/metrics.json} and \texttt{models/personalloan/metrics.json}. These files also contain confusion matrices and results for independent synthetic and boundary datasets. The root \texttt{models/model\_metadata.json} contains a legacy model that is not to be used in place of the product metrics. The product-specific artifacts are reported in Chapter 7 and the legacy results are kept separate.

\section{Prediction Outputs and Local Explanations}
\subsection{Prediction Probability}
Let $p$ be the model's probability for approval and $\hat y$ its predicted class. The response confidence value is
\begin{equation}
  Q=\begin{cases}p,&\hat y=1,\\1-p,&\hat y=0.\end{cases}
\end{equation}
It is the model probability of the predicted class. These product models do not have a separate calibration method like isotonic regression or Platt scaling in the repository. Thus, it is not an explanation-reliability measure or a verified calibration estimate, but a prediction-confidence representation.

\subsection{SHAP Implementation}
The model registry initializes a \texttt{SHAPExplainer} with \texttt{shap}. For every XGBoost classifier, we use TreeExplainer. For each application the transformed feature vector is passed to the explainer. The implementation returns one feature record for each model feature transformed, extracts the contribution vector, maps the sign to positive, negative or neutral direction. Local contribution values are not percentages of the probability of approval, but rather values in the output space of the SHAP explainer. When displaying raw values for one-hot and engineered features, be careful that the feature name does not not match a raw applicant field.

\subsection{LIME Implementation}
The \texttt{LIMEExplainer} is a wrapper for \texttt{lime.lime\_tabular.LimeTabularExplainer}. The explainer is given the transformed training background and model feature names. For an explanation of an input, it invokes the \texttt{predict\_proba} method of the classifier and asks for the votes for the accepted class. The route uses a deterministic seed generated from the input JSON, via the SHA-256 function, before executing LIME. This is to ensure repeatable sampling of the same input in a consistent software configuration, but does not guarantee that all versions of a library or all hardware configurations will yield identical results bitwise.

\section{Explanation Validation and Trust Score}
The route performs three explanation checks and a separate Gender-sensitivity check. The following formulas describe the current project definitions and are not universal standards.

\subsection{Faithfulness}
For each method, the implementation retains up to five features with the largest absolute contributions and uses their signed terms with the existing baseline or intercept. Let $p$ be the model's approval probability and $g_m$ the resulting approximation for method $m$. SHAP contributions are reconstructed in log-odds and converted to probability using the sigmoid function; LIME uses its interpretable feature values and intercept. Local concordance and faithfulness are
\begin{equation}
 LC_m=\max(0,1-|p-g_m|),\qquad
 F=0.60\,LC_{\mathrm{SHAP}}+0.40\,LC_{\mathrm{LIME}}.
 \label{eq:faithfulness}
\end{equation}
The 60/40 weights are a project choice. The check evaluates the original input without masking or refitting. Invalid or inconsistent reconstruction information makes the check unavailable. Close reconstruction is not a probability that the explanation is correct and does not establish faithful behaviour for other inputs.

\subsection{Stability}
The implementation constructs up to 20 distinct nearby applications by changing one eligible financial input at a time and recalculating dependent features. For each valid neighbour, it compares the original and new SHAP explanations using the clipped Spearman correlation $R$ of absolute-contribution ranks and the direction agreement $D$ on the union of their top-five features. The overall score is the mean of the $n$ valid pairwise comparisons:
\begin{equation}
  S=\frac{1}{n}\sum_{i=1}^{n}\left(0.70R_i+0.30D_i\right).
\end{equation}
The check also records the lowest similarity and model decision changes. Fewer than two distinct valid processed neighbours make the check unavailable. The score describes the tested small changes and does not guarantee stability for all applicants or changes.

\subsection{Consistency}
Let $A_5$ and $B_5$ be the sets of five features with the largest absolute contributions under SHAP and LIME. Let $D_c$ be the fraction of shared features with matching nonzero contribution directions. The consistency score is
\begin{equation}
  C=0.60\frac{|A_5\cap B_5|}{|A_5\cup B_5|}+0.40D_c.
  \label{eq:consistency}
\end{equation}
If there is no comparable shared direction, $D_c$ is zero. The code does not compare raw contribution magnitudes across the two methods because their output scales differ.

\subsection{Individual Fairness Check}
For a normalized applicant profile $x$, the code obtains approval probabilities $p(x)$ and $p(x^{(g)})$, where $x^{(g)}$ has the gender field flipped and other fields held fixed. The score is
\begin{equation}
  R=\max\left(0,\min\left(1,1-|p(x)-p(x^{(g)})|\right)\right).
  \label{eq:counterfactual}
\end{equation}
This is an individual model-sensitivity check for one binary attribute. If it cannot be performed, the routine returns zero rather than a perfect score; that value does not distinguish failure from maximum sensitivity. It is not an estimate of group approval-rate parity, equal opportunity, causal fairness, or legal compliance.

\subsection{Combined Trust Score}
For new assessments, the Trust Score is the arithmetic mean of the three explanation diagnostics, each on $[0,1]$:
\begin{equation}
  T=\frac{F+S+C}{3}.
  \label{eq:trust-score}
\end{equation}
The equal weights are defined in \texttt{backend/app/trust/trust\_score.py}; the Gender-flip score is reported separately and does not contribute to $T$. Previously saved assessments retain their recorded scores and weights. This experimental summary has no external calibration study or validated decision threshold.

\section{Application Implementation}
\subsection{Backend and API Communication}
The single evaluation is done by the \texttt{POST /predict} endpoint, which is mounted under \texttt{/api/v1/loan}. Batch endpoints take CSV uploads or JSON rows. Other endpoints return, list, and delete evaluation or batch history, return sample CSVs and presets, and create a PDF report. Pydantic schemas are used to define application fields, optional product fields, explanation structures, solvency data, and response objects.

The frontend makes a JSON request for single predictions and PDF creation, and multipart form data for CSV uploads. The answer contains prediction, approval probability, predicted-class confidence representation, model version, SHAP feature array, LIME feature array, Trust Score components, solvency summary, and dossier ID. Each output is individually accessible to the UI and stored history thanks to the API's structured response.

\subsection{Frontend and Application Interfaces}
The React application uses a screen state to select a portal mode, loan product, application form, and result view. The bank-officer and applicant screens are two different presentations of the same evaluation results. The dashboard displays prediction information, SHAP/LIME summaries, validation metrics, submitted form values, and a loan-specific what-if affordability calculator. The what if calculator recalculates EMI and FOIR using fixed illustrative rates and does not call the model when slider changes. It should not therefore be interpreted as a prediction of approval.

\paragraph*{Application Interfaces}
Figures~\ref{fig:loan-selection-screenshot}--\ref{fig:pdf-screenshot} show product selection, the application form, officer and applicant results, saved history and the PDF report. These views connect the submitted application with its assessment and stored record.

\clearpage
\begingroup
\setlength{\intextsep}{5pt}
\setlength{\textfloatsep}{8pt}
\setlength{\floatsep}{8pt}
\renewcommand{\topfraction}{1.0}
\renewcommand{\bottomfraction}{1.0}
\renewcommand{\textfraction}{0.0}
\reportfigure{figures/screenshots/loan_selection.png}{0.99}{7.0}{Loan Selection Interface}{Loan selection for the home-loan and personal-loan products.}{fig:loan-selection-screenshot}

\reportfigure{figures/screenshots/loan_application_form.png}{0.99}{12.5}{Loan Application Interface}{Home-loan application form.}{fig:application-screenshot}
\clearpage

\reportfigure{figures/screenshots/bank_officer_results.png}{0.99}{9.5}{Bank-Officer Results Interface}{Bank-officer summary and review guidance.}{fig:officer-screenshot}

\reportfigure{figures/screenshots/applicant_results.png}{0.99}{10.2}{Applicant Results Interface}{Applicant assessment and affordability summary.}{fig:applicant-screenshot}
\clearpage

\reportfigure{figures/screenshots/application_history.png}{0.99}{9.8}{Application History Interface}{Application history with saved assessments and report actions.}{fig:history-screenshot}

\reportfigure{figures/screenshots/pdf_report_preview.png}{0.99}{11.5}{Generated PDF Interface}{Generated loan assessment report.}{fig:pdf-screenshot}
\clearpage
\endgroup

\subsection{Database and Application History}
Evaluation and batch information are stored in the models \texttt{PredictionLog} and \texttt{BatchUpload} of SQLAlchemy. Local SQLite is always initialized as a fallback, a PostgreSQL URL can be provided via \texttt{DATABASE\_URL}. If the primary engine fails, the request session tries the secondary engine (SQLite). The data in the database contains JSON data for raw inputs, full responses, explanations, and bad batch data. This means that the project can be run without a remote database, but that local SQLite and cloud PostgreSQL are viable options (not a replicated multi-database consistency design).

\subsection{PDF Report Generation}
The reporting module generates an A4 PDF containing application reference, applicant and facility summary, decision rationale, solvency information, trust and validation values, and next steps using ReportLab. The PDF endpoint is used to accept applicant data and prediction response and return a PDF response. The report is generated on demand -- no cryptographic document hash or signed officer approval record is created by the code.

\clearpage
\section{End-to-End System Workflow}
\noindent\begin{minipage}{\textwidth}
\textbf{Algorithm 2: Online loan evaluation}\par\smallskip
\begin{lstlisting}
Input: LoanApplicationRequest x
Output: prediction response y
1. Select product pipeline M from x.loan_type
2. Normalize product-specific values and construct one-row data frame
3. Derive EMI, FOIR, and (for home loans) LTV features
4. Transform the row with the product's fitted preprocessor
5. Obtain model class y_hat and approval probability p
6. Generate SHAP and LIME local explanations
7. Calculate faithfulness, stability, consistency, and gender-flip scores
8. Calculate weighted Trust Score and predicted-class probability value
9. Generate product-prefixed dossier ID
10. Persist record when database session is available; return response
\end{lstlisting}
\end{minipage}\par\medskip

This sequence illustrates the behaviour in the API route. It also demonstrates the need for an explanation-quality score to be separate from the prediction probability and the deterministic FOIR summary.

\par\medskip\noindent
This chapter has followed the implementation from data preparation, feature engineering, training, local explanation, validation and the online assessment path. It has also outlined the presentation, storage and export of the response, separate from the financial summary and Trust Score.
